% Part:sets-functions-relations
% Chapter: size-of-sets
% Section: pairing

\documentclass[../../../include/open-logic-section]{subfiles}

\begin{document}

\olfileid{sfr}{siz}{pai}
\olsection{जोडीकरण फलने आणि संकेतांक}

\begin{explain}
कँटर यांच्या नागमोडी पद्धतीमुळे $\Nat^n$ ची गणनीयता चित्रातून
स्पष्ट दिसते. पण $\Nat^2$ दाखवणाऱ्या आपल्या सरणीवर लक्ष केंद्रित करू.
सरणीतील नागमोडी रेषेचा मागोवा घेऊन स्थाने मोजल्यास,
$\tuple{1,2}$ ही जोडी~$7$ या संख्येशी जोडली आहे हे तपासता येते.
पण ही संख्या अधिक थेट मोजता आली, तर सोयीचे होईल. म्हणजे, नागमोडी
प्रगणनाचे \emph{व्युत्क्रम} $g\colon \Nat^2 \to \Nat$ हाताशी असावे,
आणि त्यासाठी
\[
g(\tuple{0,0}) = 0, \;
g(\tuple{0,1}) = 1, \;
g(\tuple{1,0}) = 2, \; \dots,  
g(\tuple{1,2}) = 7, \; \dots
\]
असे असावे. त्यामुळे $\tuple{n, m}$ ही जोडी आपल्या प्रगणनात नेमकी
कोणत्या स्थानी येईल ते मोजता येईल.

खरे तर दोन निरीक्षणांच्या साहाय्याने $g$ ची थेट व्याख्या करता येते.
पहिले निरीक्षण: $n$ व्या ओळीतील आणि $m$ व्या स्तंभातील मूल्य~$v$
असेल, तर $(n+1)$ व्या ओळीतील आणि $(m-1)$ व्या स्तंभातील मूल्य
$v + 1$ असते. दुसरे निरीक्षण: आपल्या प्रगणनाची पहिली ओळ $0$, $1$,
$3$, $6$, इत्यादींपासून सुरू होणाऱ्या त्रिकोणी संख्यांची आहे.
$k$ वी त्रिकोणी संख्या म्हणजे $< k$ ही अट पूर्ण करणाऱ्या नैसर्गिक संख्यांची
बेरीज; ती $k(k+1)/2$ अशी मोजता येते. ही दोन निरीक्षणे एकत्र करून
पुढील फलन विचारात घ्या:
\[
  g(n,m) = \frac{(n+m+1)(n+m)}{2} + n
\]
$g(\tuple{n, m})$ ऐवजी आपण बहुधा फक्त $g(n, m)$ असे लिहितो, कारण
ते पाहायला सोपे असते. या सूत्रानुसार आधी $(n+m)^\text{वी}$ त्रिकोणी
संख्या काढायची आणि नंतर तिच्यात $n$ मिळवायचा. त्यामुळे सरणी नेमकी
आपल्याला हवी तशी भरते. विशेषतः, $\tuple{1, 2}$ ही जोडी
$\frac{4 \times 3}{2} + 1 = 7$ कडे पाठवली जाते.

हे फलन $g$ म्हणजे जोड्यांच्या संचाच्या एका प्रगणनाचे \emph{व्युत्क्रम}
आहे. अशा फलनांना \emph{जोडीकरण फलने} म्हणतात.
\end{explain}

\begin{defn}[जोडीकरण फलन]
  $f\colon A \times B \to \Nat$ यातील $f$ हे एकास-एक असेल, तर ते
  \emph{अंकगणिती जोडीकरण फलन} असते. $f$ हे $A \times B$ चे
  \emph{सांकेतीकरण करते} आणि $f(x,y)$ हा $\tuple{x,y}$ चा
  \emph{संकेतांक} असतो असेही आपण म्हणतो.
\end{defn}

\begin{explain}
जोडीकरण फलनांनी, उदाहरणार्थ, नैसर्गिक संख्यांच्या जोड्यांचे
सांकेतीकरण करता येते; म्हणजेच घटकांची प्रत्येक \emph{जोडी} एका
\emph{एकाच} संख्येने दर्शवता येते. जोडीकरण फलनाचे व्युत्क्रम वापरून
संख्येचे \emph{विसांकेतीकरण} करता येते, म्हणजे ती संख्या कोणती जोडी
दर्शवते हे शोधता येते.
\end{explain}

\begin{prob}
सर्व ऋणेतर परिमेय संख्यांच्या संचाचे एक प्रगणन द्या.
\end{prob}

\begin{prob}
$\Rat$ हा !!{enumerable} आहे, हे दाखवा. कोणतीही परिमेय संख्या
$z/m$ या अपूर्णांकाच्या रूपात लिहिता येते, जेथे $z \in \Int$ आणि
$m \in \Nat^+$ असतात, हे आठवा.
\end{prob}

\begin{prob}
$\Bin^*$ चे एक प्रगणन परिभाषित करा.
\end{prob}

\begin{prob}
प्रास्ताविक तर्कशास्त्राच्या अभ्यासक्रमातून आठवा की प्रत्येक संभाव्य
सत्यताकोष्टक एक सत्यता-फलन व्यक्त करतो. दुसऱ्या शब्दांत, सत्यता-फलने
म्हणजे कोणत्या तरी~$k$ साठी $\Bin^k \to \Bin$ अशी सर्व फलने आहेत.
सर्व सत्यता-फलनांचा संच गणनीय आहे, हे सिद्ध करा.
\end{prob}

\begin{prob}
कोणत्याही अनंत !!{enumerable} संचाच्या सर्व सांत उपसंचांचा संच
!!{enumerable} आहे, हे दाखवा.
\end{prob}

\begin{prob}
$\Nat$ च्या एखाद्या सांत उपसंचाचा पूरक असलेल्या $\Nat$ च्या उपसंचाला
\emph{सांत-पूरक} म्हणतात. म्हणजे, $A \subseteq \Nat$ हा सांत-पूरक
असतो तेव्हा आणि केवळ तेव्हाच, जेव्हा $\Nat\setminus A$ हा सांत असतो.
%READERNOTE{OLSIZ-002: गोठवलेल्या इंग्रजी वाक्यात “complement of a finite set Nat” अशी संबंधदर्शक शब्दांशिवाय अपूर्ण रचना आहे. लगेचचे सूत्र अभिप्रेत अर्थ ठरवते: Nat मधील एखाद्या सांत उपसंचाचा Nat-सापेक्ष पूरक. मराठी मजकुरात हा संबंध स्पष्ट केला आहे.}
$\Nat$ चे सांत आणि सांत-पूरक उपसंच हेच ज्याचे !!{element}s आहेत,
असा संच $I$ घ्या. $I$ हा !!{enumerable} आहे, हे दाखवा.
\end{prob}

\begin{prob}
!!{enumerable} इतक्या !!{enumerable} संचांचा संयोग !!{enumerable} असतो,
हे दाखवा. म्हणजे, $A_1$, $A_2$, \dots{} हे संच असतील आणि प्रत्येक
$A_i$ हा !!{enumerable} असेल, तर त्या सर्वांचा संयोग
$\bigcup_{i=1}^\infty A_i$ हाही !!{enumerable} असतो. [टीप: हे कठीण आहे!]
\end{prob}

\begin{prob}
$f \colon A \times B \to \Nat$ हे कोणतेही जोडीकरण फलन असू द्या.
हे $f$ आच्छादकही असेल, तर त्याचे व्युत्क्रम $A \times B$ चे प्रगणन
आहे, हे दाखवा. सर्वसाधारण बाबतीत व्युत्क्रम केवळ फलनाच्या व्याप्तीवर
परिभाषित असते; यावरून जोड्यांचा हा संच गणनीय आहे, हे सिद्ध करा.
  %READERNOTE{MRSIZ-002: गोठवलेल्या इंग्रजी सरावात कोणत्याही एकास-एक जोडीकरण फलनाचे व्युत्क्रम थेट प्रगणन असल्याचा दावा आहे. फलन आच्छादक नसेल, तर त्याचे व्युत्क्रम सर्व नैसर्गिक संख्यांवर परिभाषित नसते. मराठी सरावात आच्छादक बाब आणि सामान्य अंशतः-व्युत्क्रम बाब वेगळ्या केल्या आहेत.}
\end{prob}

\begin{prob}
$\Nat^3$ चे सांकेतीकरण करणारे फलन स्पष्टपणे द्या.
\end{prob}

\end{document}
